\chapter{FUNDAMENTAÇÃO TEÓRICA}

\orientacao{Esta seção apresenta os conceitos, teorias e referenciais que sustentam sua pesquisa. Organize em seções temáticas (não uma referência por parágrafo). Cada seção deve construir a base teórica necessária para entender o problema e a metodologia. Cite autores relevantes com \texttt{\textbackslash cite\{\}}.}

\section{\placeholder{Título da primeira seção teórica}}

\orientacao{Exemplo de títulos: ``Conceitos fundamentais de [tema]'', ``Estado da arte em [área]'', ``Frameworks e modelos relacionados''.}

\placeholder{Apresente os conceitos centrais do seu tema. Defina termos-chave. Explique teorias ou modelos que fundamentam sua abordagem.}

\placeholder{Continue desenvolvendo o tópico com citações de autores consagrados na área. Relacione os conceitos entre si.}

\section{\placeholder{Título da segunda seção teórica}}

\placeholder{Aborde um segundo eixo teórico relevante para o seu trabalho. Por exemplo: tecnologias, metodologias de avaliação, contexto regulatório, etc.}

\placeholder{Conecte esta seção com a anterior e com o problema de pesquisa apresentado na introdução.}

\section{\placeholder{Título da terceira seção teórica — opcional}}

\placeholder{Inclua quantas seções forem necessárias para cobrir a base teórica do seu estudo. Remova seções não utilizadas.}

\orientacao{Encerre a fundamentação teórica indicando a lacuna ou oportunidade de pesquisa que seu trabalho endereça — isso prepara o leitor para a seção de Trabalhos Relacionados e para a Metodologia.}

\placeholder{Parágrafo de transição: ``Apesar dos avanços descritos, ainda são escassos estudos que [lacuna]. Nesse sentido, este trabalho propõe [sua contribuição].''}
